\begin{abstract}
We present Rubin Data Preview 2 (DP2), the second data preview from the National Science Foundation--Department of Energy Vera C. Rubin Observatory, comprising raw and calibrated single-epoch images, coadds, difference images, detection catalogs, and ancillary data products.
DP2 is derived from observations acquired by the LSST Science Camera (LSSTCam) on the Simonyi Survey Telescope at the Summit Facility on Cerro Pach\'on, Chile, primarily during the on-sky commissioning campaign between \dptwostartdate\ and \svcampaignenddate, supplemented by observations taken between \startoperations and \dptwoenddate  that overlap the commissioning footprint.
The DP2 footprint comprises the Science Validation wide-area survey, five Deep Drilling Fields, and a number of targeted small-field regions, including Trifid and Lagoon, Prawn, M49, and New Horizons, all observed as part of the Rubin First Look campaign.
Each field was imaged in up to six broad photometric bands, $ugrizy$, and coadded to produce deep imaging covering an estimated \coaddedarea.
The addition of single-visit-only areas expands the total DP2 footprint to an estimated \dptwoarea, with coverage in at least one filter.
The median per-visit PSF FWHM across the wide-area survey ranges from \zmedianpsffwhmwide in the $z$ band to \grmedianpsffwhmwide in $g$ and $r$ bands, with the sharpest fields reaching $\sim$ \bestimagequality.
COSMOS, the deepest field, reaches estimated coadded 5$\sigma$ depths of $u$=\umaglim, $g$=\gmaglim, $r$=\rmaglim, $i$=\imaglim, $z$=\zmaglim, $y$=\ymaglim.
DP2 contains \nobjects distinct astrophysical objects, of which \nextendedobjects are extended in at least one band in coadds, and \nsolarsystemobjects solar system objects, of which \nnewasteroiddiscoveries are new discoveries.
DP2 is approximately \sizeinbytes in size and is available to Vera C. Rubin Observatory data rights holders via the Rubin Science Platform, a cloud-based environment for the analysis of petascale astronomical data.
Based on a roughly five-month primary observing baseline and covering only part of the eventual LSST footprint, DP2's area, depth, and multiband coverage nonetheless support a broad range of early science investigations ahead of LSST Data Release 1.
\end{abstract}
